\section{A spanning theorem on the principal charts}

% BEGIN B76S1
\paragraph{B76S1. Choosing the harmonic family in the kernel.}
For the two specified row weights take their finite-dimensional
complex space $V$ of harmonic row-highest polynomials. Its total
degree is fixed by the weights. In the stable range $p+q\le n$
it is nonzero. For each embedding use the flag minors of the first
kind in one set of $p$ column positions and those of the second
kind in a disjoint set of $q$ positions. Take their product with
the exponents prescribed by the two weights. Each factor has
a nonzero diagonal monomial. Addition of an earlier row inside
a flag minor leaves its determinant unchanged, and diagonal row
scaling gives the prescribed weight. Disjoint column positions
kill every mixed contraction. This gives the original minor
polynomial, including the constant when both weights vanish.

Equip $V$ with the Fock inner product, in which monomials are
orthogonal with squared norm $m!$ for multi-index $m$.
Simultaneous unitary column changes preserve the row equations
and the contractions, hence $V$. They preserve the inner product:
identify degree $d$ polynomials with symmetric tensors of degree
$d$ in the variable space carrying its Hermitian metric. Up to a
single degree-dependent factor the induced metric gives the
displayed monomial norms, as seen by counting the permutations
of a multiset of indices. This applies to both the defining
unitary action and its dual. Every such change therefore takes an
orthonormal basis $P_1,\ldots,P_m$ of $V$ to another.
For any two linear evaluations $e,e'$ on $V$ the scalar
$\sum_i e(P_i)\overline{e'(P_i)}$ is invariant under that change
of basis by unitarity. Apply this to evaluation at two independent
arguments. We obtain precisely the nonzero invariant family
required in the heat derivative and the real reproducing formula.
% END B76S1

% BEGIN B76S2
\paragraph{B76S2. Spanning over the given coefficient field.}
Fix a pinned lower row and its principal charts of C76P1. Take the
identity cover of ICL11, the canonical connection of CNC15, a
dominant pair $(a,n)$ (in particular $n$ is even) and
$n\ge4[E:\mathbb Q]r+4$. For every original section $f$ whose
readings are bounded on these charts,
\[
 f\in\operatorname{span}_{c.F}
      \mathrm{shortGen}(X_o,c,\nabla,1,0,a,n,d).
\]
In particular C76P2 gives this conclusion when boundedness is
known in the distinguished uniformization. All weights and
coefficient embeddings are those of the given row.

Use one adapted marking and the reference Hermitian space of the
principal chart construction, with its level tie, letter dictionary
and stabilizer calculation from C76P3--C76P7. Its moved component lattices $L_k$
have disjoint orbits and trivial stabilizer multipliers. On each
representative transport $C_k^{-1}\mathrm{rd}(f)_k$ by the inverse
change of coordinates. Principal congruence and the trivial component
character give its stabilizer invariance with the prescribed Borel
weight. For any transported lattice $\gamma L_k$ extend this
function by multiplying by $\mu(\gamma,L_k)^{-1}$ and applying
$\gamma^{-1}$ to the frame. If two representatives give the same
lattice their indices agree and their quotient is in its stabilizer;
the cocycle identity then makes both extensions equal. Set the form
zero on the remaining orbits. Holomorphy is preserved on each
component, and the finite maximum of the component bounds bounds
the new form because the action is unitary and preserves the
Gram minors in the norm.

Take the finite system at the auxiliary integer $2M$, where
$M=l.M$. This lies in the required level ideal and forces the
definite lattices to be trivial at the required primes. L76F6
supplies the nine system fields. The doubled product, its
Heisenberg change and E76Q1--E76Q3 supply the comparison laws on
these very tests. Apply I75K3 with B76S1 and the real integral
of K74A4. One finite expansion holds for the extended form for
all $L_k$ and frames. Since $D.\mathfrak m=\mathrm{Mfr}(l)$,
S76G2 identifies each test in it with the required geometric
test. To choose $N$, triviality of $\Lambda$ at $2M$ gives a
positive integer $t$ prime to $2M$ with $tS_0\subseteq\Lambda$
and $t\Lambda\subseteq S_0$, where $S_0=\mathcal O_E^n$. In
particular $N=Mt$ is odd, divisible by $M$, and
$NS_0\subseteq\Lambda\cap S_0$. This is an allowed choice for
each definite lattice. Expand
each harmonic polynomial by ICL15 in polynomials integral on that
lattice. These choices precede evaluation in any component.
A76T2 gives the original sections, and S76G4 replaces each
analytic term by its section reading divided by $C_ku_{\Lambda,N}$.
Multiplying back by $C_k$ leaves coefficients independent of $k$.
Outside the frame domain readings are zero by definition.

We have obtained a finite complex dependence of the readings. If
$f$ lay outside the $c.F$-span of the finitely many constructed
sections, a basis of their span together with $f$ would be independent
over $c.F$. Its readings would be independent over $\mathbb C$ by
ICL8, a contradiction. This proves (ICL.11) on the chosen original
pinned data. Selection of such data for every family with the
required eventual bound is a further quantifier in the whole
obligation.
% END B76S2
